\definecolor{myred}{RGB}{200, 30, 30}

\begin{figure}[h!]
\centering

% ABOVE
\begin{subfigure}[c]{0.32\textwidth}
\centering
\begin{tikzpicture}[scale=0.7]
  \coordinate (qIn)  at (-2.0, 1.4);
  \coordinate (k1)   at (-1.1, 0.5);
  \coordinate (gISR) at (-0.4, 1.6);
  \coordinate (qbIn) at (-2.0, -1.4);
  \coordinate (V1)   at (-0.5, 0);
  \coordinate (V2)   at (1.5, 0);
  \coordinate (k2)   at (2.5, 1);
  \coordinate (gln2)   at (3, 1.4);
  \coordinate (gFSR) at (3, 0.5);
  \coordinate (gln1)  at (3, -1.4);

  \draw[fermion] (qIn)  -- (V1);
  \draw[fermion] (V1)   -- (qbIn);
  \draw[gluon]   (k1)   -- (gISR);
  \draw[gluon]   (V1)   -- (V2);
  \draw[scalararrow, myred]   (V2)   -- (gln2);
  \draw[scalararrow, myred]   (gln1)  -- (V2);
  \draw[gluon]   (k2)   -- (gFSR);

  \node[left]       at (qIn)  {\(q\)};
  \node[left]       at (qbIn) {\(\bar{q}\)};
  \node[above left] at (gISR) {\(g\)};
  \node[right, myred]      at (gln1) {\(\overline{\tilde{t}}\)};
  \node[below]      at (gFSR) {\(g\)};
  \node[right, myred]      at (gln2) {\(\tilde{t}\)};
  \node[above]      at ( 0.5, 0.0) {\(g\)};
\end{tikzpicture}
\end{subfigure}
\hfill
\begin{subfigure}[c]{0.32\textwidth}
\centering
\begin{tikzpicture}[scale=0.7]
  \coordinate (qIn)  at (-2.0, 1.4);
  \coordinate (k1)   at (-1.1, 1.4);
  \coordinate (gISR) at (-0.4, 2);
  \coordinate (qbIn) at (-2.0, -1.4);
  \coordinate (V1)   at (+0.5, 1.4);
  \coordinate (V2)   at (+0.5, -1.4);
  \coordinate (k2)   at (2.1, 1.4);
  \coordinate (gln1)   at (3, 1.4);
  \coordinate (gFSR) at (3, 0.5);
  \coordinate (gln2)  at (3, -1.4);

  \draw[fermion] (qIn)  -- (V1);
  \draw[fermion] (V2)  -- (qbIn);
  \draw[gluon, myred] (V1)   -- (V2);
  \draw[myline, myred]    (V1)   -- (V2);
  \draw[gluon]   (k1)   -- (gISR);
  \draw[scalararrow, myred]   (V1)   -- (gln1);
  \draw[scalararrow, myred]   (gln2)   -- (V2);
  \draw[gluon]   (k2)   -- (gFSR);

  \node[left]       at (qIn)  {\(q\)};
  \node[left]       at (qbIn) {\(\bar{q}\)};
  \node[above left] at (gISR) {\(g\)};
  \node[right, myred]      at (gln1) {\(\tilde{t}\)};
  \node[below]      at (gFSR) {\(g\)};
  \node[right, myred]      at (gln2) {\(\overline{\tilde{t}}\)};
  \node[right, myred]      at (-0.2, 0.0) {\(\tilde{g}\)};
\end{tikzpicture}
\end{subfigure}
\hfill
\begin{subfigure}[c]{0.32\textwidth}
\centering
\begin{tikzpicture}[scale=0.7]
  \coordinate (qIn)  at (-2.0, 1.4);
  \coordinate (k1)   at (-0.9, 0.6);
  \coordinate (gISR) at (-0.4, 1.6);
  \coordinate (qbIn) at (-2.0, -1.4);
  \coordinate (V1)   at (+0.5, 0);
  \coordinate (k2)   at (2.5, 1);
  \coordinate (gln2)   at (3, 1.4);
  \coordinate (gFSR) at (3, 0.5);
  \coordinate (gln1)  at (3, -1.4);

  \draw[gluon] (qIn)  -- (V1);
  \draw[gluon] (V1)  -- (qbIn);
  \draw[gluon]   (k1)   -- (gISR);
  \draw[scalararrow, myred]   (gln1)   -- (V1);
  \draw[scalararrow, myred]   (V1)   -- (gln2);
  \draw[gluon]   (k2)   -- (gFSR);

  \node[left]       at (qIn)  {\(g\)};
  \node[left]       at (qbIn) {\(g\)};
  \node[above left] at (gISR) {\(g\)};
  \node[right, myred]      at (gln1) {\(\overline{\tilde{t}}\)};
  \node[below]      at (gFSR) {\(g\)};
  \node[right, myred]      at (gln2) {\(\tilde{t}\)};
  \node[right]      at (-0.2, 0.0) {\(\tilde{q}\)};
\end{tikzpicture}
\end{subfigure}

\par\medskip
%DOWN
\begin{subfigure}[c]{0.32\textwidth}
\centering
\begin{tikzpicture}[scale=0.7]
  \coordinate (qIn)  at (-2.0, 1.4);
  \coordinate (k1)   at (-1.1, 0.5);
  \coordinate (gISR) at (-0.4, 1.6);
  \coordinate (qbIn) at (-2.0, -1.4);
  \coordinate (V1)   at (-0.5, 0);
  \coordinate (V2)   at (1.5, 0);
  \coordinate (k2)   at (2.5, 1);
  \coordinate (gln2)   at (3, 1.4);
  \coordinate (gFSR) at (3, 0.5);
  \coordinate (gln1)  at (3, -1.4);

  \draw[gluon] (qIn)  -- (V1);
  \draw[gluon] (V1)   -- (qbIn);
  \draw[gluon]   (k1)   -- (gISR);
  \draw[gluon]   (V1)   -- (V2);
  \draw[scalararrow, myred]   (V2)   -- (gln2);
  \draw[scalararrow, myred]   (gln1)   -- (V2);
  \draw[gluon]   (k2)   -- (gFSR);

  \node[left]       at (qIn)  {\(g\)};
  \node[left]       at (qbIn) {\(g\)};
  \node[above left] at (gISR) {\(g\)};
  \node[right, myred]      at (gln1) {\(\overline{\tilde{t}}\)};
  \node[below]      at (gFSR) {\(g\)};
  \node[right, myred]      at (gln2) {\(\tilde{t}\)};
  \node[above]      at ( 0.5, 0.0) {\(g\)};
\end{tikzpicture}
\end{subfigure}
\hfill
\begin{subfigure}[c]{0.32\textwidth}
\centering
\begin{tikzpicture}[scale=0.7]
  \coordinate (qIn)  at (-2.0, 1.4);
  \coordinate (k1)   at (-1.1, 1.4);
  \coordinate (gISR) at (-0.4, 2);
  \coordinate (qbIn) at (-2.0, -1.4);
  \coordinate (V1)   at (+0.5, 1.4);
  \coordinate (V2)   at (+0.5, -1.4);
  \coordinate (k2)   at (2.1, 1.4);
  \coordinate (gln1)   at (3, 1.4);
  \coordinate (gFSR) at (3, 0.5);
  \coordinate (gln2)  at (3, -1.4);

  \draw[gluon] (qIn)  -- (V1);
  \draw[gluon] (V2)  -- (qbIn);
  \draw[scalararrow, myred] (V2)   -- (V1);
  \draw[gluon]   (k1)   -- (gISR);
  \draw[scalararrow, myred]   (V1)   -- (gln1);
  \draw[scalararrow, myred]   (gln2)   -- (V2);
  \draw[gluon]   (k2)   -- (gFSR);

  \node[left]       at (qIn)  {\(g\)};
  \node[left]       at (qbIn) {\(g\)};
  \node[above left] at (gISR) {\(g\)};
  \node[right, myred]      at (gln1) {\(\tilde{t}\)};
  \node[below]      at (gFSR) {\(g\)};
  \node[right, myred]      at (gln2) {\(\overline{\tilde{t}}\)};
  \node[right, myred]      at (-0.2, 0.0) {\(\tilde{t}\)};
\end{tikzpicture}
\end{subfigure}
\hfill
\begin{subfigure}[c]{0.32\textwidth}
\centering
\begin{tikzpicture}[scale=0.7]
  \coordinate (qIn)  at (-2.0, 1.4);
  \coordinate (k1)   at (-1.1, 1.4);
  \coordinate (gISR) at (-0.4, 2);
  \coordinate (qbIn) at (-2.0, -1.4);
  \coordinate (V1)   at (+0.5, 1.4);
  \coordinate (V2)   at (+0.5, -1.4);
  \coordinate (V3)   at (+1.75, 0);
  \coordinate (k2)   at (2.5, 0.9);
  \coordinate (gln2)   at (3, 1.4);
  \coordinate (gFSR) at (3, 0.5);
  \coordinate (gln1)  at (3, -1.4);

  \draw[gluon] (qIn)  -- (V1);
  \draw[gluon] (V2)  -- (qbIn);
  \draw[scalararrow, myred] (V1)   -- (V2);
  \draw[gluon]   (k1)   -- (gISR);
  \draw[scalararrow, myred]   (V1)   -- (V3);
  \draw[scalararrow, myred]   (V3)   -- (gln1);
  \draw[scalararrow, myred]   (V2)   -- (V3);
  \draw[scalararrow, myred]   (V3)   -- (gln2);
  \draw[gluon]   (k2)   -- (gFSR);

  \node[left]       at (qIn)  {\(g\)};
  \node[left]       at (qbIn) {\(g\)};
  \node[above left] at (gISR) {\(g\)};
  \node[right, myred]      at (gln1) {\(\tilde{t}\)};
  \node[below]      at (gFSR) {\(g\)};
  \node[right, myred]      at (gln2) {\(\overline{\tilde{t}}\)};
  \node[right, myred]      at (-0.2, 0.0) {\(\tilde{t}\)};
\end{tikzpicture}
\end{subfigure}

\caption{Representative leading-order Feynman diagrams for the production of stop pairs; each one shown with ISR and FSR. Adapted from~\cite{10.21468/SciPostPhysCore.7.4.072}.}
\label{fig:hscp_prod_stop}
\end{figure}